\label{sec:introduction}

Every member of the House of Representatives maintains district offices---physical
offices located within the district where constituents can meet with staff, receive
assistance with federal agency matters, and engage with their representative's
legislative work. The number of district offices per member ranges from one to seven,
depending on district geography and appropriated office budgets. For the average
constituent, the district office is the primary point of physical contact with
their congressional representative; the Capitol Hill office is accessible only by
traveling to Washington.

The geographic distance between a constituent and their representative's district
office is therefore a direct measure of representational access---not a proxy
for it. A constituent who lives 15 minutes from their representative's office can
drop in without losing a workday; a constituent who lives 90 minutes away faces
a significant barrier to in-person engagement. This access asymmetry compounds
across the approximately 200 million Americans who contact their congressional
representatives annually for constituent services.

Despite the intuitive importance of this access dimension, we are not aware of any
systematic empirical study linking district geography to constituent travel distance
at national scale. The redistricting literature has extensively studied compactness
as a legal and aesthetic criterion \citep{polsby1991,young1988}, documented its
correlation with partisan gerrymandering \citep{chen2013}, and used it as a predictor
of various electoral outcomes \citep{levin2019}. But the direct welfare interpretation
of compactness---that compact districts reduce the travel burden on constituents---has
not been quantified.

This paper fills that gap. Our analysis has four components:

\begin{enumerate}
\item We collect the district office addresses for all 435 members of the 118th Congress
  from \texttt{house.gov} (the official U.S. House website), identifying the primary
  district office address for each member.

\item We compute population-weighted centroid coordinates for all 74,134 census tracts
  in the contiguous 48 states using 2020 P.L. 94-171 population data.

\item We route each tract centroid to the relevant representative's primary district
  office address using the OSRM routing API with the OpenStreetMap North America
  road network, recording the estimated driving time in minutes.

\item We aggregate tract-level driving times to district-level means (weighted by
  tract population) and analyze the correlation with district Polsby-Popper compactness.
\end{enumerate}

The key finding---a correlation of $r = -0.67$ between Polsby-Popper and mean driving
time---is robust across different compactness metrics (Reock: $r = -0.59$, Schwartzberg:
$r = -0.61$) and is substantially larger in the subset of districts with documented
partisan gerrymandering. The \textsc{bisect} comparison shows a 14-minute reduction
in mean constituent driving time under algorithmic plans relative to enacted gerrymanders.

\subsection{Relationship to the O-Series Framework}

This paper is the fourth in the O-series (O.0--O.4), which investigates redistricting
effects on democratic outcomes. The O.0 framework paper establishes the four outcome
dimensions and identification strategy. The O.4 analysis is distinctive within the
series in two respects: it measures an outcome (travel time) that is directly computable
from geography rather than requiring electoral data, and it provides a causal mechanism
(the ``neck'' geometry) rather than relying on the quasi-experimental design used in
O.1--O.3. This makes the O.4 findings less subject to the confounding concerns that
motivate the DiD and IV designs in O.1--O.3.

\subsection{Policy Significance}

The finding has immediate policy applications. Several state redistricting statutes
require that districts be ``compact,'' but courts have struggled to give this
requirement operational content \citep{levin2019}. The constituent distance metric
provides a concrete, welfare-relevant definition: a compact district is one whose
constituents can efficiently access their representative. This definition is
non-partisan, directly measurable, and explains why compactness matters as a
criterion beyond its role in preventing partisan manipulation.
